% The shape of the learned encoder: blocks halving the grid, channels doubling up to dim_h, the
% mean over the last map, the head; and the receptive field of the blocks.
\begin{tikzpicture}[fig]
% (a) blocks
\node[lab, anchor=west] at (-0.2,2.25) {(a) a network of 4 blocks, width 32, $\dim h=256$};
\foreach \k/\s/\c/\lab in {0/1.6/cP2!15/{input $n$}, 1/1.2/cP0!15/{$n/2$, 32}, 2/0.9/cP0!30/{$n/4$, 64},
                         3/0.65/cP0!45/{$n/8$, 128}, 4/0.45/cP0!65/{$n/16$, 256}} {
  \pgfmathsetmacro{\x}{2.0*\k}
  \fill[\c, draw=cInk2] (\x,1.0-\s/2) rectangle ++(\s,\s);
  \node[small] at ({\x+\s/2},-0.05) {\lab};
}
\foreach \k in {0,...,3} \draw[->] ({2.0*\k+1.65},1.0) -- ({2.0*\k+1.95},1.0);
\node[small] at (4.3,-0.55) {each block: grid $/2$, channels $\times2$ up to $\dim h$};
\draw[->] (8.5,1.0) -- (9.0,1.0) node[right, small] {mean};
\node[draw, fill=cP1!20, minimum width=0.25cm, minimum height=1.1cm, label={[small]above:$h$}] (h) at (10.0,1.0) {};
\draw[->] (10.15,1.0) -- (10.7,1.0);
\node[draw, rounded corners, small, align=center] (g) at (11.6,1.0) {$g$: hidden,\\ projection};
\draw[->] (g) -- (12.9,1.0) node[right, small] {$z$};
\node[small, text width=3.1cm] at (11.6,-0.05) {training only; loss at temperature $T$};
% (b) receptive field
\begin{scope}[yshift=-3.0cm]
  \node[lab, anchor=west] at (-0.2,1.9) {(b) receptive field and frame};
  \foreach \r/\b in {0.13/1, 0.34/2, 0.77/3} \draw[cP0, thick] ({1.5-\r},{0.6-\r}) rectangle ({1.5+\r},{0.6+\r});
  \node[small, anchor=west] at (3.0,1.4) {1 block: 7 points};
  \node[small, anchor=west] at (3.0,0.9) {2 blocks: 19};
  \node[small, anchor=west] at (3.0,0.4) {3 blocks: 43; 4 blocks: 91};
  \node[note, text width=6.6cm, align=left, anchor=west] at (6.2,0.6)
    {\textbf{Checked before training:} last map of at least 2 points per side (blocks $\le\log_2(n/2)$);
     channels never decrease (width $\le\dim h$); projection $\le\dim h$ and $\le$ hidden; values in
     their lists or ranges. Shown: receptive field against the frame side $n$, number of weights,
     rough memory of a step.};
\end{scope}
\end{tikzpicture}
